%%% FIELD revealed-input-statement
At \(N=4\), \(m=1\), and \(\alpha=2/5\), let \(Q^{\mathrm{id}}\) be the revealed-input channel on \(\mathcal Z=\mathcal X\), defined by \(Q^{\mathrm{id}}(z\mid x)=\mathbf 1\{z=x\}\). Then
\[
 V_1(Q^{\mathrm{id}})\ge \frac{637}{240}.
\]
For every finite-output common channel \(Q\), whether private or not, \(V_1(Q)\ge V_1(Q^{\mathrm{id}})\). Consequently \(V_1(Q)\ge637/240\) uniformly over every local-privacy class.
%%% FIELD revealed-input-proof
Write an orbit as \(n=(n_{00},n_{01},n_{10},n_{11})\), enumerate \(\mathcal O_4\) lexicographically, and put
\[
 P^{\mathrm{id}}_n(h)=P_{n,1,Q^{\mathrm{id}}}(h),
 \qquad
 h=(h_{00},h_{01},h_{10},h_{11}).
\]
The orbit-kernel definition applies because complete randomization assigns exactly one of the four units to treatment and the identity channel reports each realized input \((W_i,Y_i(W_i))\) without error.

We first give an exact dual certificate rather than assuming the displayed lower bound. Define nonnegative arrays \(\pi=(\pi_n)\) and \(\eta=(\eta_n)\) by the following table; every unlisted entry is zero.
\[
\begin{array}{c|c|c@{\qquad}c|c|c}
n&\pi_n&\eta_n&n&\pi_n&\eta_n\\ \hline
(0,0,4,0)&0&1/16 &(0,2,0,2)&0&1/4\\
(0,3,0,1)&1/4&5/12 &(0,4,0,0)&0&1/16\\
(1,0,0,3)&0&1/6 &(1,0,1,2)&0&1/4\\
(1,0,2,1)&0&3/8 &(1,0,3,0)&1/4&5/12\\
(1,1,0,2)&1/4&1/2 &(1,1,2,0)&1/4&0\\
(1,2,0,1)&0&3/8 &(2,0,0,2)&0&5/8\\
(2,0,1,1)&0&1/2 &(2,0,2,0)&0&1/4\\
(2,1,0,1)&0&1/4 &(3,0,0,1)&0&1/6
\end{array}
\]
Direct addition gives
\[
 \sum_{n\in\mathcal O_4}\pi_n=1,
 \qquad
 \sum_{n\in\mathcal O_4}\eta_n=\frac{14}{3}.
\]
For a report histogram \(h\) and an action \(t\in\mathcal T_4\), define
\[
 b_h=\sum_n\pi_nP^{\mathrm{id}}_n(h),
 \qquad
 e_{h,t}=\sum_{n:\,\tau(n)=t}\eta_nP^{\mathrm{id}}_n(h).
\]
The identity kernel is especially transparent. If the unique treated unit has type \((a,b)\), its report is \((1,b)\), while every other unit of type \((a',b')\) reports \((0,a')\). Therefore the orbit-kernel coefficient formula reduces to
\[
 P^{\mathrm{id}}_n(h)
 =\frac14\sum_{a,b\in\{0,1\}}n_{ab}
 \mathbf 1\!\left\{
 h=\left(
 n_{00}+n_{01}-\mathbf 1\{a=0\},
 n_{10}+n_{11}-\mathbf 1\{a=1\},
 \mathbf 1\{b=0\},
 \mathbf 1\{b=1\}
 \right)\right\}.
\]
Substitution of the displayed \(\pi\) and \(\eta\) into this formula gives exactly four nonzero positive parts:
\[
\begin{array}{c|c|c}
h&t&[e_{h,t}-b_h]_+\\ \hline
(0,3,1,0)&-3/4&1/24\\
(0,3,1,0)&-1/2&1/32\\
(3,0,0,1)& 1/2&1/32\\
(3,0,0,1)& 3/4&1/24
\end{array}
\]
All other positive parts vanish. Hence
\[
 \sum_{h,t}[e_{h,t}-b_h]_+
 =\frac1{24}+\frac1{32}+\frac1{32}+\frac1{24}
 =\frac7{48}.
\]

For completeness, let \((c,\ell)\) be any feasible point of the fixed-channel confidence program. Weighting its loss constraints by \(\pi_n\), using \(\sum_n\pi_n=1\), and then weighting its coverage constraints by \(\eta_n\) gives
\[
 \ell\ge \sum_{h,t}b_hc_{h,t},
 \qquad
 \sum_{h,t}e_{h,t}c_{h,t}
 \ge (1-\alpha)\sum_n\eta_n.
\]
Because \(0\le c_{h,t}\le1\), subtraction yields the weak-duality inequality
\[
\begin{aligned}
 \ell
 &\ge \sum_{h,t}e_{h,t}c_{h,t}
      -\sum_{h,t}(e_{h,t}-b_h)c_{h,t}\\
 &\ge (1-\alpha)\sum_n\eta_n
      -\sum_{h,t}[e_{h,t}-b_h]_+\\
 &=\frac35\frac{14}{3}-\frac7{48}
 =\frac{637}{240}.
\end{aligned}
\]
Taking the infimum over feasible \((c,\ell)\) proves \(V_1(Q^{\mathrm{id}})\ge637/240\).

Finally, fix any finite-output common channel \(Q\). From the identity transcript \(X_i=(W_i,Y_i(W_i))\), independently draw \(Z_i\sim Q(\cdot\mid X_i)\) and then apply any randomized confidence rule designed for \(Z\). By conditional product randomization, the resulting confidence set has, for every fixed schedule, exactly the same distribution, coverage, and expected cardinality as in the experiment using \(Q\) directly. Thus every feasible rule for \(Q\) induces a feasible rule with the same worst-schedule loss for \(Q^{\mathrm{id}}\). Taking infima gives
\[
 V_1(Q^{\mathrm{id}})\le V_1(Q),
\]
which proves the claimed uniform bound.
%%% FIELD strengthened-theorem-statement
For \(N=4\) and \(\alpha=2/5\), every global optimizer of \(L^\star_{4,2/5}(\lambda)\) uses the balanced allocation \(m=2\) for every \(\lambda\ge1000\). At \(\lambda=1000\),
\[
 \frac{129113957}{50000000}
 \le L^\star_{4,2/5}(1000)
 \le \frac{64557127}{25000000}.
\]
Moreover,
\[
 [999.9,1000.1]\cap\mathcal B_{4,2/5}=\varnothing,
\]
so no balanced randomized response after a nontrivial binary partition is globally optimal anywhere in that interval. Separately, at \(N=4\), \(m=2\), \(\lambda=3\), and \(\alpha=1/10\), the mapping \(Q\mapsto-V_2(Q)\) is not convex, so the confidence objective does not satisfy the convex-utility premise of the generic extreme-point shortcut.
%%% FIELD strengthened-theorem-proof
For \(j\in\{1,2,3\}\), define the allocation-specific value
\[
 U_j(\lambda)=\inf_{Q\in\mathcal D_\lambda}V_j(Q).
\]
This is an auxiliary notation, while the causal design target remains
\[
 L^\star_{4,2/5}(\lambda)=\min_{j\in\{1,2,3\}}U_j(\lambda).
\]

The exact four-unit certificate at \(\lambda=1000\) supplies a feasible four-output channel and confidence rule at \(m=2\) with loss at most
\[
 U_2(1000)\le \frac{64557127}{25000000}.
\]
If \(\lambda\ge1000\), then every \(Q\in\mathcal D_{1000}\) also belongs to \(\mathcal D_\lambda\): indeed,
\[
 Q(z\mid x)\le1000Q(z\mid x')
 \le\lambda Q(z\mid x')
\]
for every \(x,x'\in\mathcal X\) and \(z\), where nonnegativity of channel probabilities and \(\lambda\ge1000\) justify the second inequality. The same certified channel and the same confidence rule therefore remain feasible, with their schedulewise laws and loss unchanged. Consequently
\[
 U_2(\lambda)\le \frac{64557127}{25000000}
 \qquad(\lambda\ge1000).
\]

The revealed-input lower-bound lemma applies to every channel, without a privacy restriction, and hence gives
\[
 U_1(\lambda)\ge\frac{637}{240}
 \qquad(\lambda>1).
\]
Allocation reflection, applied within the same common-channel class \(\mathcal D_\lambda\), gives
\[
 U_3(\lambda)=U_1(\lambda)\ge\frac{637}{240}.
\]
The separation is strict because
\[
 \frac{637}{240}-\frac{64557127}{25000000}
 =\frac{5391119}{75000000}>0.
\]
Thus, for every \(\lambda\ge1000\), neither \(m=1\) nor \(m=3\) can minimize the three allocation-specific values, whereas \(m=2\) has a strictly smaller feasible upper bound. Polynomial attainment supplies an optimizer for each allocation-specific channel game, and the allocation set is finite. It follows that every global optimizer uses \(m=2\) throughout the asserted ray.

At \(\lambda=1000\), the exact simultaneous polynomial-dual certificate also gives
\[
 U_2(1000)\ge\frac{129113957}{50000000}.
\]
Because the preceding strict comparison excludes \(m=1,3\), one has \(L^\star_{4,2/5}(1000)=U_2(1000)\), which proves the displayed bracket.

We next verify the interval-specific nonbinary conclusion. For \(r\)-ary randomized response, direct subtraction of a row at \(\lambda\) and at \(\lambda_0\) gives
\[
 \operatorname{TV}\!\left\{Q_{r,\lambda}(\cdot\mid x),
 Q_{r,\lambda_0}(\cdot\mid x)\right\}
 =\frac{(r-1)|\lambda-\lambda_0|}
 {(\lambda+r-1)(\lambda_0+r-1)}.
\]
Set \(\lambda_0=1000\). When \(|\lambda-1000|\le1/10\), the sum of the product-experiment total-variation bounds from the experiment-value modulus lemma for the binary and four-output channels is at most
\[
 4\left\{
 \frac{1/10}{1000^2}+
 \frac{3/10}{1000^2}
 \right\}
 =\frac{16}{10\cdot1000^2}.
\]
Here \(4=N\), and each displayed denominator is a valid lower bound because \(\lambda+1\ge1000.9\), \(\lambda_0+1=1001\), \(\lambda+3\ge1002.9\), and \(\lambda_0+3=1003\). Applying the value modulus separately to the two channels and adding gives a total value perturbation at most
\[
 (2N+1)(1+\alpha^{-1})\frac{16}{10\cdot1000^2}
 =9\left(1+\frac{1}{2/5}\right)
   \frac{16}{10\cdot1000^2}
 =\frac{63}{1250000}.
\]
At \(\lambda_0=1000\), the exact certificate gives a feasible four-output value
\[
 U=\frac{25823}{10000}
\]
and, for every nontrivial binary partition \(S\), the lower bound
\[
 V_2(R_{S,1000})>D,
 \qquad D=\frac{258239}{100000}.
\]
Their certified gap dominates the perturbation because
\[
 D-U-\frac{63}{1250000}
 =\frac{9}{100000}-\frac{63}{1250000}
 =\frac{99}{2500000}>0.
\]
If \(\delta_2(\lambda)\) and \(\delta_4(\lambda)\) denote the two value-modulus bounds, respectively, the preceding calculation states \(\delta_2(\lambda)+\delta_4(\lambda)\le63/1250000\). Thus, for every \(\lambda\in[999.9,1000.1]\) and every nontrivial \(S\),
\[
\begin{aligned}
 V_2(R_{S,\lambda})-V_2(Q_{4,\lambda})
 &\ge V_2(R_{S,1000})-V_2(Q_{4,1000})
       -\delta_2(\lambda)-\delta_4(\lambda)\\
 &>D-U-\frac{63}{1250000}>0.
\end{aligned}
\]
Therefore the continued four-output channel has strictly smaller fixed-channel value than every \(R_{S,\lambda}\). Since the unrestricted value is no greater than that four-output value, the definition of the balanced-binary gap implies
\[
 \Delta_{4,2/5}(\lambda)>0,
\]
and hence \([999.9,1000.1]\cap\mathcal B_{4,2/5}=\varnothing\).

Finally, take \(N=4\), \(m=2\), \(\lambda=3\), and \(\alpha=1/10\). The exact certificate provides a binary channel \(Q_b\), a four-output channel \(Q_f\), and their tagged midpoint \(Q_{\mathrm{mix}}=\tfrac12Q_b\oplus\tfrac12Q_f\), satisfying
\[
 V_2(Q_b)>\frac{347}{50},
 \qquad
 V_2(Q_f)>\frac{737}{100},
 \qquad
 V_2(Q_{\mathrm{mix}})\le\frac{711}{100}.
\]
It follows that
\[
 V_2(Q_{\mathrm{mix}})
 <\frac{1431}{200}
 <\frac{V_2(Q_b)+V_2(Q_f)}2.
\]
This violates the midpoint inequality required for concavity of \(V_2\). Equivalently, \(Q\mapsto-V_2(Q)\) is not convex, proving the final clause.
